% Auto-generated by scripts/bench_signal_heal_latency.py — do not hand-edit numbers.
\paragraph{Transfer inference latency.}
Table~\ref{tab:transfer-latency} reports wall-clock inference on transfer holdout batches (device=\texttt{cpu}, batch=64, warmup=20, timed=100; includes prolonged residual). N2N uses the wavetable \texttt{n2n\_corrupt\_corrupt} checkpoint zero-shot when present; domain-trained N2N was not run. These timings fill the prior transfer gap (wavetable tables had latency; transfer had $R$ only).

\begin{table}[t]
  \centering
  \caption{Transfer-domain inference latency (ms/batch; higher is slower). Batch=64; \texttt{---} = method unavailable.}
  \label{tab:transfer-latency}
  \setlength{\tabcolsep}{3.2pt}
  \begin{tabular}{@{}lrrrrrr@{}}
    \toprule
    Method & CWRU & MFPT & MIT-BIH & PTB-XL & CNC & PMU \\
    \midrule
    no-bake & 0.19 & 0.13 & 0.22 & 0.19 & 0.29 & 0.20 \\
    DualCosine & 0.56 & 0.58 & 0.70 & 0.60 & 0.65 & 0.63 \\
    Ours (\texttt{hybrid}) & 38.43 & 40.25 & 6.62 & 6.82 & 26.25 & 1.64 \\
    N2N (WT$\to$dom.) & 3.85 & 3.48 & 3.86 & 3.80 & 3.14 & 3.16 \\
    \bottomrule
  \end{tabular}
\end{table}

\begin{figure}[t]
  \centering
  \includegraphics[width=\columnwidth]{figures/fig_signal_heal_transfer_latency.png}
  \caption{Transfer-domain inference latency (ms/batch). N2N is wavetable-trained zero-shot when available.}
  \label{fig:signal-heal-transfer-latency}
\end{figure}

